\documentclass[a4paper]{article}
\usepackage[utf8]{inputenc}
\usepackage[natbib,sorting=none]{biblatex}
\usepackage{graphicx}
\usepackage{acronym}
\usepackage{indentfirst}
\usepackage{requirements} %custom package
\usepackage{hyphenat} %stops word separation on new line
\usepackage{fancyhdr}
\usepackage{enumitem}
\usepackage{listings}
\usepackage{xcolor}

\definecolor{commentgreen}{RGB}{2,112,10}
\definecolor{eminence}{RGB}{108,48,130}
\definecolor{weborange}{RGB}{255,165,0}
\definecolor{frenchplum}{RGB}{129,20,83}

\addbibresource{references.bib}

\begin{document}
\setlist[enumerate]{label=\alph*),leftmargin=0.4cm, topsep=0cm}
\setlength{}{}
\pagenumbering{none}
\begin{figure}[!h]
	\centering
	\includegraphics[width=80mm]{images/poli-logo.png}
\end{figure}
\hfill
\begin{center}
    \fontsize{18px}{6mm}\selectfont \textsc{\textbf{Software Engineering II project}}
\end{center}
\begin{center}
    \fontsize{12px}{4mm}\selectfont \textsc{Academic Year: 2018/2019}
\end{center}
\hfill
\hfill
\begin{figure}[!h]
	\centering
	\includegraphics[width=120mm]{images/trackme-logo.png}
\end{figure}
\hfill
\hfill
\begin{center}
    \fontsize{22px}{8mm}\selectfont \textsc{\textbf{Requirement Analysis and\\ Specification Document}}
\end{center}
\begin{center}
    \fontsize{14px}{4mm}\selectfont \textsc{{Draft Version 0.1 - 11/10/2018}}
\end{center}
\hfill
\hfill
\begin{center}
\fontsize{14px}{4mm}\selectfont \textsc{\textit{Davide Rutigliano -  903616}}
\end{center}

\begin{center}
\fontsize{14px}{4mm}\selectfont \textsc{\textit{Claudio Ferrante - 903417\\}}
\end{center}

\begin{center}
\fontsize{14px}{4mm}\selectfont \textsc{\textit{Davide Matta - 903616}}
\end{center}
\pagenumbering{roman}
\tableofcontents
\newpage
\addcontentsline{toc}{subsection}{Use Cases Summary}
\renewcommand\listtablename{Use Cases Summary}
\listoftables
\newpage
\addcontentsline{toc}{subsection}{UML Diagrams List}
\renewcommand\listfigurename{UML Diagrams List}
\listoffigures
\newpage
\rfoot{\includegraphics[width=30mm]{images/trackme-logo-mini.png}}
%%%%%%%%%%%%%%%%%%%%%%%%%%%%%%%%%%%%%%%%%%%%%%%%%%%%%%%%%%%%%%
\newpage
\pagestyle{fancy}
\pagenumbering{arabic}
\section{Introduction}

    \subsection{Purpose}
    This Requirement Analysis and Specification Document addresses the description in terms of functional and non-functional requirements of \textit{TrackMe}, a new health monitoring application.
    
    This document is aimed at giving an explanation both of the problem and of the proposed software-based solution by pointing out different features of the application. Moreover, this specifications also addresses some relevant QoS characteristics that the system should guarantee.
    
    Thus, the analysis is primarily focused on requirements elicitation and use cases description through well known standards for specification documents:
        \begin{itemize}
            \item[\textit{1)}] \textit{UML} for Use Case, Sequence and Activity Diagrams \cite{rumbaugh2004unified};
            \item[\textit{2)}] \textit{Alloy} for formal analysis \cite{jackson2006software}.
        \end{itemize}
        
    This document is mostly directed and likely necessary to software designers and developers interested in the deployment of the proposed system.
    
    \subsection{Scope and Audience}
        \subsubsection{Description of the given problem}
        
        \subsubsection{Description of the proposed system}
        \textit{TrackMe} is an ease of use health monitoring application which offers different services for both young and old people who needs to keep track of their personal data in order to to keep their health safe.
        
        This proposed system provides a very efficient data acquisition service named \textit{Data4Help}. This facility gives the possibility to verified \textit{third-party} registered onto the service to request health status information about subscribed customers. The application may also allow \textit{individual} users to connect \textit{external devices} such as smart-watches or similar, to perform a more detailed data acquisition and monitoring.
        
        In addition, the application provides also a fully automated SOS service for elderly people: \textit{Automated-SOS} uses \textit{Data4Help} features to monitor the health status and, when such parameters are below certain thresholds, sends the location of the customer to an ambulance.
    
        Moreover, through the \textit{Track4Run} service, \textit{TrackMe} allows customers to create a new run, a competition in which other users can either participate as \textit{athletes} or watch as \textit{spectators}. Additionally, this service may exploits other service functionality for keeping track of athletes progresses and for helping them to prevent major accident happen during the run. 
        \subsubsection{Goals}
        
        \begin{description}
            \item[G.01] Goal
        \end{description}
    
    \subsection{Document Overview}
        This part of the document is intended to provide both an overview of the problem and an idea of the proposed solution.
        
        Next section is aimed at giving a more detailed description of the proposed system, by meaning of the application point of view. Additionally, addresses users characteristic, dependencies and constraint.
        
        Third chapter shows a detailed analysis of \textit{TrackMe} in terms both functional and non-functional requirements of the system. Furthermore, takes into account definition of all use cases, their description through UML diagrams and their mapping on requirements.
        
        Fourth section illustrates \textit{Alloy} models including purpose, proof and explanation. In addition, describes worlds obtained by running them.
        
        Last part is about effort spent in writing this document in terms of hours of work.
        
        There is also a list of tables and a list of figures representing respectively use cases and diagrams in order to help the reader to understand them and navigate the specification.
        
    \subsection{Definition, Acronyms and  Abbreviations}
            \subsubsection{Keywords}
            The key words “MUST”, “MUST NOT”, “REQUIRED”, “SHALL”, “SHALL NOT”, “SHOULD”, “SHOULD NOT”, “RECOMMENDED”, “MAY”, and “OPTIONAL” in this document are to be interpreted as described in RFC 2119 \cite{bradner1997key}.
            \subsubsection{Definition of Terms}
            This document uses several terms that might be more loosely used elsewhere. These terms are defined here as they will be used later on in this document.
                \begin{description}
                    \item[\textbf{User}] a general definition of a customer registered into the application
                    
                    \item[\textbf{Individual}] a single person \textit{TrackMe User}, identified either by SSN or by FC number
                    
                    \item[\textbf{Third-Party}] a company \textit{TrackMe User}, identified by VAT number
                    
                    \item[\textbf{External Device}] an external device such as a smart-watch or similar devices
                    
                    \item[\textbf{Run}] a competition organized and managed by an \textit{Organizer} to which \textit{Athletes} and \textit{Spectators} can enroll as participant or watchers
                    
                    \item[\textbf{Organizer}] either an \textit{Individual} or a \textit{Third-Party} that uses \textit{Track4Run} service to create a new \textit{Run}
                    
                    \item[\textbf{Athlete}] an \textit{Individual} enrolled to an already existing \textit{Run}
                    
                    \item[\textbf{Spectator}] either an \textit{Individual} or a \textit{Third-Party} that is watching an already existing \textit{Run}
                    
                    \item[\textbf{TrackMe}] the \textit{"system to be"}
                    
                    \item[\textbf{Data4Help}] a data monitoring service provided by \textit{TrackMe}
                    
                    \item[\textbf{Automated-SOS}] an SOS service built on top of \textit{Data4Help}
                    
                    \item[\textbf{Track4Run}] run management service offered by \textit{TrackMe} application
                    
                    \item[\textbf{Credential}] as used in this document, is a combination of both username and password used by an \textit{User} to authenticate him/herself during the Log-in phase
                \end{description}
                
            \subsubsection{Abbreviations}
            \begin{itemize}
                \item[-] Abbreviation 
            \end{itemize}
            
            \subsubsection{Acronyms}
            \begin{acronym}
                \acro{RASD}{Requirement Analysis and Specification Document}
                \acro{UML}{Unified Modelling Language}
                \acro{QoS}{Quality of Service}
                \acro{SSN}{Social Security Number}
                \acro{FC}{Fiscal Code}
                \acro{VAT}{Value Added Tax}
                \acro{BT}{Bluetooth}
                \acro{NFC}{Near Field Communication}
                \acro{IU}{Individual User}
                \acro{TPU}{Third Party User}
            \end{acronym}
            
    \subsection{References}
        \printbibliography[heading=none]
%%%%%%%%%%%%%%%%%%%%%%%%%%%%%%%%%%%%%%%%%%%%%%%%%%%%%%%%%%%%%%
\newpage
\section{Overall Description}

    \subsection{Product Perspective}
    
    \subsection{Product Functions}
    
        \subsubsection{Data4Help}
        
        \subsubsection{Automated-SOS}
        
        \subsubsection{Track4Run}
    
    \subsection{User Characteristics}
    
        \subsubsection{Individual}
        
        \subsubsection{Third-Party}
        
    \subsection{Assumptions, Dependencies and Constraints}
        \begin{description}
            \item[A.01] Assumption
        \end{description}
        
        \begin{description}
            \item[D.01] Dependencies
        \end{description}
        
        \begin{description}
            \item[C.01] Constraints
        \end{description}
%%%%%%%%%%%%%%%%%%%%%%%%%%%%%%%%%%%%%%%%%%%%%%%%%%%%%%%%%%%%%%
\newpage
\section{Specific Requirements}

    \subsection{External Interface Requirements}
        
        \subsubsection{User Interfaces}
        
        \subsubsection{Hardware Interfaces}
        
        \subsubsection{Software Interfaces}
        
        \subsubsection{Communication Interfaces}
    
    \subsection{Functional Requirements}
        
        \subsubsection{Registration}
        \textit{TrackMe} must provide a registration interface for both \textit{individual} and \textit{third-party} users.
        
        An \textit{individual} must provide his/her SSN (Social Security Number) or FC (Fiscal Code) during the registration phase in order to be uniquely identified by \textit{TrackMe} application. In this way, the user is also accepting terms of service and allows \textit{TrackMe} to acquire his/her personal data.
        %use cases table structure
        \begin{usecase}{Individual User Registration}
        \name{Individual User Registration}
        \brief{An individual sign up to the application to become an user.}
        \actor{Individual, TrackMe.}
        \pre{The individual has a valid SSN/FC and he's on the main page.}
        \bflow{The individual presses a register button.}
              {TrackMe asks if user is an Individual or a Third Party. The user can choose between these two options.}
              {The individual choose the Individual option.}
              {TrackMe asks to provide a SSN/FC (based on region) and user personal information such as name, surname, age, gender, e-mail, nationality and optional notes about his health condition.}
              {The individual fills out the requested information and presses the confirm button.}
              {TrackMe registers that user as an Individual User and adds it to the User Database.}
        \post{The individual is successfully registered as an Individual User.}
        \end{usecase}
        
        A \textit{third-party} must provide a VAT number during the registration in order to allow \textit{TrackMe} to check that the submitted VAT is registered and that is issued by an EU Member State.
        
        \begin{usecase}{Third Party User Registration}
        \name{Third Party User Registration}
        \brief{A third party sign up to the application to become an user.}
        \actor{Third Party, TrackMe.}
        \pre{The third party has a valid VAT and it's on the main page.}
        \bflow{The third party presses a register button.}
              {TrackMe asks if user is an Individual or a Third Party. The user can choose between these two options.}
              {The third party choose the Third Party option.}
              {TrackMe asks if user is an Individual or a Third Party. The user can choose between these two options.}
              {The third party fills out the requested information and presses the confirm button.}
              {TrackMe registers that user as a Third Party user and adds it to the User Database.}
        \post{The third party is successfully registered as an Third Party User.}
        \end{usecase}
        
        \subsubsection{Log-in}
        The system should provide a login interface to the end user in order sign into the application.
        
        \begin{usecase}{Individual User Login}
        \name{Individual User Login}
        \brief{An individual sign in to access the application.}
        \actor{Individual, TrackMe.}
        \pre{The individual has a username and a password and it's on the main page.}
        \bflow{The individual presses a login button.}
              {TrackMe provides two boxes where the user can write. Former is titled 'username' while latter is called 'password'.}
              {The individual writes his username on the 'username' box and his associated password on 'password' box.}
              {TrackMe checks if that username is on the User DB and its associated password.}
              {TrackMe finds that username and its associated password are valid and grants access to the user.}
        \post{The individual is successfully granted privilege to access TrackMe and it's recognized as IU.}
        \end{usecase}
        
        \begin{usecase}{Third Party User Login}
        \name{Third Party User Login}
        \brief{A third party sign in to access the application.}
        \actor{Third Party, TrackMe.}
        \pre{The third party has a username and a password and it's on the main page.}
        \bflow{The third party presses a login button.}
              {TrackMe provides two boxes where the user can write. Former is titled 'username' while latter is called 'password'.}
              {The third party writes his username on the 'username' box and his associated password on 'password' box.}
              {TrackMe checks if that username is on the User DB and its associated password.}
              {TrackMe finds that username and its associated password are valid and grants access to the user.}
        \post{The third party is successfully granted privilege to access TrackMe and it's recognized as TPU.}
        \end{usecase}
        
        \subsubsection{Connection of External Devices}
        The application may provide an interface that allows \textit{individual} to connect an external device such as smart-watch via Bluetooth or NFC.
        
        Moreover, in case such interface is present, \textit{TrackMe} should be able to interact also with the external device and to acquire data from it.
        
        \subsubsection{Data4Help Request}
        
        \subsubsection{Data4Help Request Management}
        
        \subsubsection{Automated-SOS Request}
        
        \subsubsection{Automated-SOS Data Management}
        
        \subsubsection{Automated-SOS Request Management}
        
        \subsubsection{Automated-SOS Data Synchronization}
        
        \subsubsection{Run Creation}
        
        \subsubsection{Run Enrollment}
        
        \subsubsection{Run Watching}
        
        \subsubsection{Run Management}
        
    \subsection{Performance Requirements}
    
    \subsection{Design Constraints}
    
        \subsubsection{Standards Compliance}
        
        \subsubsection{Hardware Limitations}
        
        \subsubsection{Other Constraints}
    
    \subsection{Software System Attributes}
        
        \subsubsection{Reliability}
        
        \subsubsection{Availability}
        
        \subsubsection{Security}
        
        \subsubsection{Maintainability}
        
        \subsubsection{Portability}
        
%%%%%%%%%%%%%%%%%%%%%%%%%%%%%%%%%%%%%%%%%%%%%%%%%%%%%%%%%%%%%%
\newpage
\section{Formal Analysis using Alloy}
\renewcommand*{\lstlistlistingname}{Alloy model samples}
\lstlistoflistings
%definition of alloy language for Latex
\lstdefinelanguage{Alloy}{
	morekeywords={
		module, open, as,
		private, abstract, sig, extends, in,
		lone, some, one, disj,
		fact, pred, fun, assert,
		run, check,
		for, but, exactly,
		this, not, implies, else, let,
		not, no, set, all, sum,
		iff, or, Int, and,
		none, univ, iden, \#
	},
	sensitive=true,
	morecomment=[l]{//},
	morecomment=[l]{--},
	morecomment=[s]{/*}{*/},
	morestring=[b]{"},
	literate={->}{$\rightarrow$}1
	% replacing characters can cause problems when copying from PDF to editor
}[keywords,comments,strings]
\lstset {
    language=Alloy,
    frame=tb,
    tabsize=3,
    showstringspaces=false,
    numbers=none,
    commentstyle=\color{commentgreen},
    keywordstyle=\color{blue},
    stringstyle=\color{red},
    basicstyle=\small\ttfamily, % basic font setting
    escapechar=\&,
    % keyword highlighting
    classoffset=1, % starting new class
    keywordstyle=\color{weborange},
    classoffset=0,
}
\lstinputlisting[language=Alloy, caption=Model Definition, firstline=0, lastline=60]{TrackMeAlloyModel.als}
\begin{lstlisting}
\end{lstlisting}
%%%%%%%%%%%%%%%%%%%%%%%%%%%%%%%%%%%%%%%%%%%%%%%%%%%%%%%%%%%%%%
\newpage
\section{Effort Spent}
    \begin{itemize}
        \item[-] \textbf{Davide Rutigliano:}
        
        \item[-] \textbf{Davide Matta:}
        
        \item[-] \textbf{Claudio Ferrante:}
    \end{itemize}

\end{document}
